\section{Experimental design}
\label{sec:experimental-design}

The experiments test whether the protocol does what Section~\ref{sec:method-rule} says, how much of the dependence uncertainty
the band spans, and what each component contributes. They run the engine on the committed documents of the case and on planted
dependence structures whose capital is known by simulation. Every number in Section~\ref{sec:results} comes from a script that
is committed beside the manuscript and is run from the repository root.

\subsection{The predictions}
\label{sec:design-hypotheses}

The nine predictions were set by the author for this study; the design record of the protocol states no acceptance criteria, so they
were not taken from a prior specification. They are written to \texttt{experiments/hypotheses.json}, which the experiment scripts
read. By file time they were written on 6 November 2026, before the runs of the reproduction, planted-structure and ablation
experiments (H2 to H4c, H6 and H7). Their order relative to the first run of the structure experiment (H1 and H5) is not recorded,
and the file was first versioned together with the results. Each prediction is stated with the outcome that refutes it.
\begin{description}
\item[H1, structure.] In each of the six committed aggregates the band has eight rows in the stated order, exactly one row labelled
headline capital and it is the simple sum, seven rows labelled information, one horizon and one confidence level, and an Euler
allocation that adds up to the aggregate expected shortfall to a relative error of at most $10^{-6}$. Refuted if any aggregate fails
any element.
\item[H2, reproduction.] The engine, given the marginals rebuilt from the committed quantile grids and the committed seed, returns
every row of the committed band within the larger of three combined Monte Carlo standard errors and 0.5 percent of the committed
row. Refuted if any row differs by more.
\item[H3, the band is not noise and no row speaks for it.] The width of the band, the worst case less the lower bound, is at least
100 times the largest Monte Carlo standard error of its rows, and for every row some other row differs from it by more than 10
percent. Refuted if the width is below 100 standard errors or some row lies within 10 percent of every other row.
\item[H4a, coverage by the sharp range.] Over the planted structures the interval from the lower bound to the worst case contains
the capital of the structure in all of them, after widening it by three standard errors of the capital of the structure.
\item[H4b, coverage by the interior rows.] The interval spanned by the five interior rows of the band (the square-root rule, the Gaussian copula
and the $t$ copulas with 4, 8 and 30 degrees of freedom) contains the capital of the structure in at least 75 percent of the
structures. Refuted if fewer are covered. ``The band'' here is the base band, computed with the base matrix, which is the stricter
reading; Section~\ref{sec:results-planted} also reports the reading with the interior rows of both matrices, so the refutation does
not depend on the reading.
\item[H4c, a single figure is not reliable.] The Gaussian copula row alone lies within 5 percent of the capital of the structure in
at most half of the structures. Refuted if more than half are within 5 percent.
\item[H5, invariance of the bounds.] The lower bound and the worst case are equal in the base band and in the band under the
market-credit correlation of 0.8 (relative difference below $10^{-9}$), while the Gaussian copula row moves by more than 10
percent.
\item[H6, the common horizon and level matter.] Reading one marginal, the market marginal, at 99 percent in place of 99.9 percent
changes the simple sum by at least 10 percent, and the consistency test returns red on that band.
\item[H7, each row feeds a test.] Removing the $t$ rows leaves the copula-choice test unrunnable and removing the stressed band
leaves the correlation-stress test unrunnable; with the Gaussian row alone, at most half of the nine tests can be run.
\end{description}
H4a is a consequence of Proposition~\ref{prop:coverage}; its outcome verifies the implementation and carries no information about
the structures.

\subsection{Data and provenance}
\label{sec:design-data}

The data are the committed documents of case C18 of Contraste, the aggregation of capital under macroeconomic scenarios, for
Banco Andino. The bank is synthetic and its loss models are simulated, so no statement below is about the capital of any
actual bank. There are four scenario documents, the baseline, the adverse scenario of the European Banking Authority 2025, the
severely adverse scenario of the Federal Reserve 2026 and the delayed-transition scenario of the Network for Greening the
Financial System, and a fifth document, the $t$ copula variant of the baseline, which also holds the band under the
market-credit correlation of 0.8. The five documents carry one and the same aggregation block, shown by its SHA-256 digest, and
the block of the stress is the sixth aggregate: the six aggregates reduce to $\bdDistinct$ distinct ones. The marginals come from four producers: the credit marginal from the engine's
grouped copula on six portfolio segments with 400 representative names each, the market marginal from the committed mixed-portfolio
forecast of the market case moved from 99 percent and one day to 99.9 percent and one year by the normal quantile ratio and the
square root of 250, the interest rate marginal from the engine's one-factor Hull--White simulation on the curve and the book of
the interest-rate case, and the operational marginal from the Fourier aggregation of a Poisson--generalised Pareto compound
whose parameters are those of the operational-risk case scaled by assets.

The documents commit the expected loss, the value at risk, the expected shortfall and the capital of each marginal and, for each
marginal, its quantile function on a grid of $\bdGridPoints$ probabilities, the standard normal distribution function at the points from
$-4.9$ to $4.9$ in steps of $0.005$, which reaches probabilities of $5\times10^{-7}$ at each end; the simulated samples are not
committed. The experiments rebuild each marginal by linear interpolation of the committed grid, which is the same rule that
produced the committed band, and the quantile function is clamped at the ends of the grid. The market marginal has losses below
zero, a market gain, down to CLP \bdMarketGridMin{} billion at the lower end of its grid. Table~\ref{tab:marginals} gives the
four marginals.

\subsection{Planted dependence structures}
\label{sec:design-planted}

To know what the band contains, the capital of the aggregate must be known under a dependence structure that is not read from the
band. Ten structures are planted on the four marginals of Banco Andino. Four are Gaussian copulas: with the base matrix, with
the matrix of the market-credit stress, with a crisis matrix in which the three financial risks have a correlation of 0.8 with
each other and 0.4 with operational risk, and with independent risks. Three are $t$ copulas: with 3 degrees of freedom and the
base matrix, with 2.5 degrees of freedom and the base matrix, and with 3 degrees of freedom and the crisis matrix; the degrees
of freedom are below the set of the band. One is a mixture of two Gaussian regimes in which 5 percent of the draws come from the
crisis matrix and the others from the base matrix. Two are copulas of families that are not in the band: a reflected Clayton copula,
which has upper tail dependence as a loss tail needs, with the parameter $\theta=\bdClaytonTheta$ that gives the mean Kendall rank
correlation $\bar\tau=\bdTauBar$ of the base matrix, and a Gumbel copula with $\theta=\bdGumbelTheta$ and the same mean. The Clayton and Gumbel
samplers are written for this study, because the engine does not simulate them, and each is checked against its closed
form: the Laplace transform of the stable variable used by the Gumbel sampler, and the rank correlation of each copula
(simulated $\bdClaytonTauSim$ against $\bdClaytonTauExact$ for the Clayton copula at $\theta=2$, and $\bdGumbelTauSim$ against
$\bdGumbelTauExact$ for the Gumbel copula at $\theta=1.67$; the difference of 0.006 for the Gumbel check is twice the standard
error of Kendall's coefficient on 50,000 pairs). The ten structures are the choice of the author: they span matrices, degrees of
freedom and families inside and outside the band, and they are not a sample from a population of dependence structures of banks.

\subsection{Methods compared}
\label{sec:design-methods}

Every method is run by the product on the same marginals, and none is taken from its publication. The rows of the band are the
simple sum, the square-root rule, the Gaussian copula \citep{li2000,rosenberg2006}, the $t$ copula \citep{demarta2005} at
4, 8 and 30 degrees of freedom, and the sharp bounds by the rearrangement algorithm \citep{puccetti2012,epr2013}. Two further
methods are single-figure estimators that a bank could use in place of the band, given the most favourable input: data from the
planted structure itself. The hybrid value at risk of Rosenberg and Schuermann \citep{rosenberg2006} combines the quantile of each
marginal with the empirical Pearson correlation of a sample of the structure, $\sqrt{v^\top R_{\mathrm{emp}}v}$ with $v_k$ the
sample quantile of marginal $k$, less the expected losses. The calibrated Gaussian copula takes the rank correlation of a sample of
the structure, converts it to a correlation by $\rho=\sin(\pi\tau/2)$, the inverse of Equation~\ref{eq:kendall}, and simulates the
Gaussian copula with that matrix. Both are calibrated on a sample of 200,000 draws of the structure (the rank correlation on its first 50,000), kept apart from
the draws that define its capital. They represent a bank that observes the joint tail with a long sample and which can estimate a matrix perfectly, so
their errors are expected to be smaller than those of a bank.

\subsection{Metrics}
\label{sec:design-metrics}

For a planted structure $s$ the capital $b^\ast(s)$ is the mean of $R_s=20$ independent estimates of Equation~\ref{eq:coupling-capital}
on $10^6$ draws each, with the standard error $\mathrm{se}^\ast(s)=\mathrm{sd}/\sqrt{R_s}$. The error of a single figure $b_j$ is
\begin{equation}
e_j(s)=\frac{b_j}{b^\ast(s)}-1 ,
\label{eq:relative-error}
\end{equation}
and a figure is within the tolerance $\delta$ of the structure when $\lvert e_j(s)\rvert\le\delta$, with $\delta=0.05$ unless
stated. An interval $[\ell,u]$ covers the structure when, after widening by three standard errors of the capital of the structure,
\begin{equation}
\ell-3\,\mathrm{se}^\ast(s)\le b^\ast(s)\le u+3\,\mathrm{se}^\ast(s),
\label{eq:coverage-indicator}
\end{equation}
and the coverage of a set of $M$ structures is the share of those covered. The width of an interval relative to the capital is
$(u-\ell)/b^\ast(s)$; it is the price of coverage, because an interval that is wide enough covers everything. The reproduction of
a committed row uses the combined standard error $\sqrt{\mathrm{se}_{\mathrm{rebuilt}}^2+\mathrm{se}_{\mathrm{committed}}^2}$, and the
width of the band in standard errors is $(b_{\mathrm{worst}}-b_{\mathrm{lower}})/\max_j\mathrm{se}_j$, the maximum taken over the
simulated rows. The difference between the rows is measured by the largest relative difference from a row to another row,
$\max_{l\ne j}\lvert b_l/b_j-1\rvert$. The relative error is signed, so that an underestimate, the dangerous direction for a
capital figure, is told apart from an overestimate, and the median and the maximum of its absolute value summarise it over the
structures. The widening by three standard errors is a common convention that makes the indicator insensitive to the
estimation error of the capital of the structure; it favours an interval by less than the tolerance of 5 percent, since the
standard error of every capital is under one percent of it.

\subsection{Repetitions, seeds and statistics}
\label{sec:design-statistics}

The committed band was produced with the seed 42, $10^6$ draws and $10^5$ cells for the bounds, and the reproduction (H2) runs the
engine at that seed. The variation of the band over the seed is measured by 20 independent runs of the base band (seeds 1001 to
1020) and 10 of the stressed band (seeds 2001 to 2010), reported as the mean and the standard deviation over runs. The capital of
planted structure number $s$, counted from 0, uses the seeds $3001+100s+r$ for $r=0,\dots,19$, its calibration sample the seed
$9001+s$, and the calibrated Gaussian copula the seed $9501+s$. The Euler allocation is repeated over the 20 seeds 4001 to 4020. The
tests of H1, H5, H6 and H7 are deterministic functions of the committed documents and of the engine, and have no repetitions. The
interval of every uncertainty is the standard deviation over seeds, or three standard errors where an interval is widened. No
hypothesis test with a $p$-value is used: each prediction is a stated threshold on a measured quantity, and the table of outcomes
reports the measured value beside the threshold.

\subsection{Ablations planned}
\label{sec:design-ablation}

Four components of the protocol are removed one at a time and the loss is measured on the planted structures. The sharp bounds are
removed, leaving the five interior rows of both matrices. The stressed matrix is removed, leaving the interior rows of the base
matrix. The $t$ rows are removed, leaving the square-root rule and the Gaussian row. And each row is used alone, as a single figure,
with its coverage at the tolerance of 5 percent. The component that is not an interval, the common horizon and level, is ablated
by reading one marginal at another level (H6), and the checks are ablated by removing the rows each reads (H7). The Euler
allocation is compared with the allocation by stand-alone expected shortfall, which does not add up. Each ablation is evaluated on
the same ten structures and the same rows, so that the variants differ only in the component removed and the capital of the
structures is not re-estimated; the interval ends are the means of the independent runs of the band, so that they do not carry
the sampling error of a single run.
